\documentclass[11pt,reqno]{amsart}
\usepackage{amsmath,amssymb,amsthm}
\usepackage[no-math]{fontspec}
\setmainfont{cmunrm.otf}[BoldFont=cmunbx.otf,ItalicFont=cmunti.otf,BoldItalicFont=cmunbi.otf]
\usepackage[english,russian]{babel}
\usepackage[margin=1in]{geometry}
\usepackage[colorlinks=true,linkcolor=blue,citecolor=blue,urlcolor=blue]{hyperref}

\newcommand{\F}{\mathbb F}
\newcommand{\Z}{\mathbb Z}
\newcommand{\G}{\mathfrak G}
\newcommand{\LL}{\mathfrak L}
\newcommand{\D}{\mathfrak D}
\newcommand{\Sm}{\mathcal S}
\newcommand{\Lh}{\widehat{\LL}}
\newcommand{\rp}[1]{\|#1\|}
\DeclareMathOperator{\spn}{span}

\theoremstyle{plain}
\newtheorem{theorem}{Теорема}
\newtheorem{lemma}[theorem]{Лемма}
\newtheorem{proposition}[theorem]{Предложение}
\newtheorem{corollary}[theorem]{Следствие}
\theoremstyle{definition}
\newtheorem{definition}[theorem]{Определение}
\newtheorem{example}{Пример}
\theoremstyle{remark}
\newtheorem{remark}[theorem]{Замечание}

\renewcommand{\andname}{и}
\renewcommand{\keywordsname}{Ключевые слова}
\renewcommand{\subjclassname}{MSC 2020}

\sloppy

\begin{document}

\title{Базис свободных обобщённых алгебр Герстенхабера}

\author{Иван Кайгородов}
\address{CMA-UBI, Университет Бейра-Интериор, Ковильян, Португалия}
\email{kaygorodov.ivan@gmail.com}

\author{Ксения Вяткина}
\address{Санкт-Петербургский государственный университет, Санкт-Петербург, Россия}
\email{vjatkina.k@gmail.com}

\thanks{Работа поддержана РНФ, проект 22-71-10001.}

\subjclass[2020]{17B63, 17B01, 17A50}
\keywords{алгебра Герстенхабера, нечётная супералгебра Пуассона, обобщённая супералгебра Пуассона, свободная алгебра, базис, супералгебра Ли}

\begin{abstract}
Обобщённые алгебры Герстенхабера --- это нечётные обобщённые супералгебры Пуассона Кантарини и Каца; алгебры Герстенхабера (нечётные супералгебры Пуассона) составляют частный случай, в котором отображение $\D(a)=\{1,a\}$ равно нулю.
Мы строим линейный базис свободной унитальной обобщённой алгебры Герстенхабера, порождённой линейно упорядоченным множеством чётных и нечётных образующих, следуя методу, применённому ранее к свободным обобщённым супералгебрам Пуассона.
В качестве приложения показано, что полилинейная часть свободной обобщённой алгебры Герстенхабера от $n$ образующих имеет размерность $n\cdot n!$.
\end{abstract}

\maketitle

\section{Введение}

Понятие скобки Пуассона восходит к работам С.~Д.~Пуассона по небесной механике начала XIX века.
С тех пор алгебры Пуассона появились во многих областях математики: пуассоновы многообразия~\cite{L1}, алгебраическая геометрия~\cite{BLLM,GK04}, некоммутативная геометрия~\cite{ber08}, операды, теория квантования~\cite{Kon03}, квантовые группы, классическая и квантовая механика.
Изучение алгебр Пуассона привело и к ряду родственных структур: обобщённые (generic) алгебры Пуассона, алгебры йордановых скобок и обобщённые алгебры Пуассона, алгебры Новикова--Пуассона, алгебры Мальцева--Пуассона--Йордана, транспонированные алгебры Пуассона, $n$-арные алгебры Пуассона, алгебры Герстенхабера~\cite{K96} и другие.

Алгебра Герстенхабера (называемая также нечётной супералгеброй Пуассона или антипуассоновой алгеброй) --- это суперкоммутативная ассоциативная супералгебра со скобкой степени $-1$, которая является скобкой супералгебры Ли относительно обращённой чётности и действует на ассоциативное умножение дифференцированиями.
Такие скобки возникают в физике, прежде всего в формализме Баталина--Вилковыского~\cite{BV81,p1,p2}, и в геометрии~\cite{p3}.
Один из важнейших чисто алгебраических источников алгебр Герстенхабера --- когомологии Хохшильда ассоциативной алгебры~\cite{Ger63}; о недавних результатах о скобках Герстенхабера на когомологиях Хохшильда см.~\cite{W21,V16,V19,L21}.
Подалгебры алгебры Герстенхабера дифференциальных операторов описаны в~\cite{B06}.

Обобщённые супералгебры Пуассона появились при изучении йордановых супералгебр.
Кантор~\cite{Kan90,Kan92} сопоставил супералгебре Пуассона йорданову супералгебру (двойник Кантора), а Кинг и МакКриммон~\cite{KM92} описали скобки, двойник Кантора которых является йордановой супералгеброй; возникающие супералгебры йордановых скобок --- это, с точностью до замены скобки, в точности обобщённые супералгебры Пуассона (см.~\cite{CK07,K17}; обобщённая версия двойника Кантора изучалась в~\cite{K10}).
В обобщённой супералгебре Пуассона правило Лейбница выполняется лишь с точностью до поправочного члена, содержащего линейное отображение $\D(a)=\{1,a\}$.
Простые линейно компактные обобщённые супералгебры Пуассона классифицированы Кантарини и Кацем~\cite{CK07}, которые затем ввели нечётные обобщённые супералгебры Пуассона и классифицировали простые линейно компактные среди них~\cite{CK10}.
Нечётные обобщённые супералгебры Пуассона с $\D=0$ --- это в точности алгебры Герстенхабера; поэтому мы называем нечётные обобщённые супералгебры Пуассона \emph{обобщёнными алгебрами Герстенхабера}.

Базисы свободных алгебр построены для свободных алгебр Ли~\cite{Hall50,shirshov,chi,Reu93}, свободных супералгебр Ли~\cite{S86,BMPZ,BKLM}, свободных супералгебр Пуассона~\cite{She93}, свободных generic-алгебр Пуассона~\cite{KSU}, свободных обобщённых супералгебр Пуассона и свободных супералгебр йордановых скобок~\cite{K17} и т.~д.
Цель настоящей заметки --- построить базис свободной унитальной обобщённой алгебры Герстенхабера, порождённой линейно упорядоченным множеством чётных и нечётных образующих (теорема~\ref{basis}), и вычислить размерность её полилинейной части (теорема~\ref{dimen}).
Мы следуем методу Ширшова~\cite{shirshov} в той форме, в какой он применён в~\cite{K17} к свободным обобщённым супералгебрам Пуассона: кандидат в базисы $U$ состоит из суперкоммутативных одночленов от элементов базиса свободной нечётной супералгебры Ли, а линейная независимость $U$ доказывается построением структуры обобщённой алгебры Герстенхабера на линейной оболочке $U$.
В нечётном случае, однако, возникает новое явление: элемент $\{1,1\}$, автоматически равный нулю во всякой обобщённой супералгебре Пуассона, является ненулевым базисным элементом свободной обобщённой алгебры Герстенхабера, так что отображение $\D$ перестаёт быть дифференцированием.
Чтобы справиться с этим, мы задаём скобку на линейной оболочке $U$ явной замкнутой формулой (раздел~\ref{sec:model}), для которой кососимметричность и правило Лейбница проверяются непосредственно, а тождество Якоби выводим из общей формулы для якобиатора (лемма~\ref{jacobiator}), которая может представлять и самостоятельный интерес.

\section{Определения и обозначения}

Всюду в статье $\F$ --- поле характеристики нуль.
Все векторные пространства и алгебры $\Z_2$-градуированы, а элементы, входящие в формулы с чётностями, предполагаются однородными.
Для однородного элемента $a$ через $|a|\in\Z_2$ обозначается его чётность, а через
$$
\rp{a}=|a|+1\in\Z_2
$$
--- его чётность относительно обращённой $\Z_2$-градуировки.

\begin{definition}\label{def}
\emph{Обобщённой алгеброй Герстенхабера} называется $\Z_2$-градуированное векторное пространство $V=V_0\oplus V_1$ с элементом $1\in V_0$ и двумя билинейными операциями $(a,b)\mapsto ab$ и $(a,b)\mapsto\{a,b\}$, такими что для всех однородных $a,b,c\in V$
\begin{align}
&|ab|=|a|+|b|,\qquad |\{a,b\}|=|a|+|b|+1;\label{id:deg}\\
&(ab)c=a(bc),\qquad ab=(-1)^{|a||b|}ba,\qquad 1a=a;\label{id:comm}\\
&\{a,b\}=-(-1)^{\rp a\rp b}\{b,a\};\label{id:skew}\\
&\{a,\{b,c\}\}=\{\{a,b\},c\}+(-1)^{\rp a\rp b}\{b,\{a,c\}\};\label{id:jacobi}\\
&\{a,bc\}=\{a,b\}c+(-1)^{\rp a|b|}b\{a,c\}+(-1)^{\rp a}\D(a)bc,\qquad\text{где } \D(a)=\{1,a\}.\label{id:leibniz}
\end{align}
Обобщённая алгебра Герстенхабера с $\D=0$ называется \emph{алгеброй Герстенхабера}.
\end{definition}

Таким образом, умножение чётно, ассоциативно и суперкоммутативно с единицей $1$, скобка нечётна, \eqref{id:skew} и \eqref{id:jacobi} --- нечётные кососимметричность и тождество Якоби, а \eqref{id:leibniz} --- обобщённое нечётное правило Лейбница.
Поскольку $\rp a\rp b=(|a|-1)(|b|-1)$ в $\Z_2$, определение~\ref{def} в точности совпадает с определением нечётной обобщённой супералгебры Пуассона из~\cite[(0.14)]{CK10}, а алгебры Герстенхабера --- в точности нечётные супералгебры Пуассона.
Основные примеры обобщённых алгебр Герстенхабера с $\D\neq0$ --- супералгебры $PO(n,n+1)$ из~\cite{CK10}.

\begin{remark}[обращение чётности]\label{rem:parity}
Обозначим через $\Pi V$ пространство $V$ с обращённой $\Z_2$-градуировкой, так что чётность элемента $a$ в $\Pi V$ равна $\rp a$.
Тождества \eqref{id:deg}, \eqref{id:skew} и \eqref{id:jacobi} означают в точности, что $(\Pi V,\{\,,\})$ --- супералгебра Ли.
Супералгебру с нечётной билинейной операцией, удовлетворяющей \eqref{id:skew} и \eqref{id:jacobi}, мы называем \emph{нечётной супералгеброй Ли}.
Обращение чётности устанавливает биекцию между нечётными супералгебрами Ли и супералгебрами Ли (ср.~\cite{Kac77} и~\cite[замечание~1.3]{CK10}); в частности, свободная нечётная супералгебра Ли, порождённая множеством $X$ чётных и нечётных элементов, --- это свободная супералгебра Ли, порождённая тем же множеством с обращёнными чётностями.
\end{remark}

\begin{remark}[отображение $\D$]\label{rem:D}
(a) Подставляя $a=1$ в \eqref{id:leibniz}, получаем
\begin{equation}\label{Dprod}
\D(bc)=\D(b)c+(-1)^{|b|}b\D(c)-\D(1)bc .
\end{equation}
Следовательно, $\D$ является нечётным дифференцированием алгебры $(V,\cdot)$ тогда и только тогда, когда $\D(1)=\{1,1\}=0$; в общем случае нечётным дифференцированием является $\D-\D(1)\,\mathrm{id}$.
В обобщённой супералгебре Пуассона (где скобка чётна) элемент $\{1,1\}$ равен нулю по кососимметричности, но \eqref{id:skew} не накладывает на $\{1,1\}$ никаких условий, поскольку $\rp 1=1$.
Мы увидим (следствие~\ref{cor:11}), что $\{1,1\}\neq0$ в свободной обобщённой алгебре Герстенхабера.

(b) Используя \eqref{id:comm} и $|\D(a)|=\rp a$, легко проверить, что \eqref{id:leibniz} равносильно следующему утверждению: для всякого $a\in V$ отображение
$$
d_a\colon b\mapsto \{a,b\}+(-1)^{\rp a}\D(a)b
$$
является супердифференцированием чётности $\rp a$ супералгебры $(V,\cdot)$.
В частности, $d_a(1)=0$, т.~е. $\{a,1\}=-(-1)^{\rp a}\D(a)$, в согласии с \eqref{id:skew}.
\end{remark}

Следующий пример показывает, что тождества определения~\ref{def} действительно не влекут $\{1,1\}=0$.

\begin{example}\label{ex:dim2}
Пусть $V=\F1\oplus\F\varepsilon$, где $|1|=0$, $|\varepsilon|=1$, умножение задано условиями $1\cdot1=1$, $1\varepsilon=\varepsilon1=\varepsilon$, $\varepsilon^2=0$, а скобка --- условиями
$$
\{1,1\}=\varepsilon,\qquad \{1,\varepsilon\}=\{\varepsilon,1\}=\{\varepsilon,\varepsilon\}=0 .
$$
Тогда $V$ --- обобщённая алгебра Герстенхабера с $\D(1)=\varepsilon\ne0$ и $\D(\varepsilon)=0$.
Действительно, \eqref{id:deg} и \eqref{id:comm} очевидны (заметим, что $\varepsilon^2=0$ вынуждено суперкоммутативностью).
Единственная ненулевая скобка базисных элементов --- $\{1,1\}$, и \eqref{id:skew} при $a=b=1$ читается как $\{1,1\}=\{1,1\}$, поскольку $\rp1=1$.
Все скобки лежат в $\F\varepsilon$ и $\{\varepsilon,V\}=\{V,\varepsilon\}=0$, поэтому обе части \eqref{id:jacobi} равны нулю.
Тождество \eqref{id:leibniz} достаточно проверить на базисных элементах.
Если $a=\varepsilon$, все члены равны нулю.
Если $a=1$ и $\varepsilon\in\{b,c\}$, каждый член содержит либо скобку с $\varepsilon$, либо произведение $\varepsilon^2$ и потому равен нулю.
Наконец, при $a=b=c=1$ тождество превращается в $\varepsilon=\varepsilon+\varepsilon-\varepsilon$.
В этой алгебре $\D(1\cdot1)=\varepsilon\ne2\varepsilon=\D(1)1+1\D(1)$, так что $\D$ не является дифференцированием, в согласии с \eqref{Dprod}.
\end{example}

\section{Свободная обобщённая алгебра Герстенхабера и множество \texorpdfstring{$U$}{U}}\label{sec:free}

Пусть $X=\{1,x_1,x_2,\dots\}$ --- конечное или счётное множество, линейно упорядоченное так, что $1<x_1<x_2<\cdots$, в котором $1$ чётно, а каждому $x_i$ приписана чётность $|x_i|\in\Z_2$.
Пусть $D_X$ --- свободная алгебра с двумя билинейными операциями $\cdot$ и $\{\,,\}$, порождённая $X$; её базис состоит из неассоциативных слов в алфавите $X$, построенных с помощью этих двух операций.
Слова однородны: $|1|=0$, $|uv|=|u|+|v|$ и $|\{u,v\}|=|u|+|v|+1$; \emph{степенью} слова называется число букв в нём.
Пусть $I$ --- идеал алгебры $D_X$, порождённый элементами
\begin{gather*}
ab-(-1)^{|a||b|}ba,\qquad (ab)c-a(bc),\qquad 1a-a,\qquad \{a,b\}+(-1)^{\rp a\rp b}\{b,a\},\\
\{a,\{b,c\}\}-\{\{a,b\},c\}-(-1)^{\rp a\rp b}\{b,\{a,c\}\},\\
\{a,bc\}-\{a,b\}c-(-1)^{\rp a|b|}b\{a,c\}-(-1)^{\rp a}\{1,a\}bc,
\end{gather*}
где $a,b,c$ пробегают однородные элементы $D_X$.
Очевидно, $\G=D_X/I$ --- свободная унитальная обобщённая алгебра Герстенхабера, порождённая $X$: любое сохраняющее чётность отображение из $\{x_1,x_2,\dots\}$ в обобщённую алгебру Герстенхабера $V$ единственным образом продолжается до гомоморфизма $\G\to V$, переводящего $1$ в $1$.

\subsection*{Хорошие слова}
Рассмотрим неассоциативные слова в алфавите $X$ относительно одной операции $\{\,,\}$.
Множество \emph{хороших слов} вместе с линейным порядком на нём определяется индукцией по степени.
Хорошие слова степени $1$ --- элементы $X$, упорядоченные как выше.
Пусть хорошие слова степени меньше $n$ определены и линейно упорядочены.
Слово $w$ степени $n$ называется хорошим, если
\begin{enumerate}
\item[(1)] $w=\{u,v\}$, где $u$ и $v$ --- хорошие слова и $u>v$;
\item[(2)] если $u=\{u_1,u_2\}$, то $u_2\le v$.
\end{enumerate}
Всякое хорошее слово степени $n$ больше всякого хорошего слова меньшей степени; два хороших слова $w=\{u_1,u_2\}$ и $w'=\{u_1',u_2'\}$ степени $n$ сравниваются лексикографически: $w>w'$, если $u_1>u_1'$ или $u_1=u_1'$ и $u_2>u_2'$.
В частности, если $w=\{u,v\}$ --- хорошее слово, то $w>u>v$.
Таким образом, хорошие слова --- это базисные коммутаторы М.~Холла~\cite{Hall50}, и они образуют множество Холла в смысле~\cite[раздел~4.1]{Reu93}.

\subsection*{Множество \texorpdfstring{$M$}{M}}
Положим
$$
M=\{\,u,\ \{v,v\}\ \mid\ u,v \text{ --- хорошие слова и } |v|=0\,\}.
$$
Порядок продолжается на $M$ по тому же правилу: элементы большей степени больше, а элементы одной степени сравниваются лексикографически по парам компонент, причём квадрат $\{v,v\}$ сравнивается как пара $(v,v)$.
Это линейный порядок на $M$ (квадрат $\{v,v\}$ никогда не является хорошим словом, так как компоненты $u_1>u_2$ хорошего слова различны), и $1$ --- наименьший элемент $M$.
Занумеруем элементы $M$: $e_1<e_2<\cdots$, так что $e_1=1$.

Пусть $\LL_X$ --- свободная алгебра с одной билинейной операцией $\{\,,\}$, порождённая $X$ (с чётностями слов, определёнными как выше), и пусть $J$ --- идеал алгебры $\LL_X$, порождённый элементами
$\{a,b\}+(-1)^{\rp a\rp b}\{b,a\}$ и $\{a,\{b,c\}\}-\{\{a,b\},c\}-(-1)^{\rp a\rp b}\{b,\{a,c\}\}$, где $a,b,c\in\LL_X$ однородны.
Фактор $\LL=\LL_X/J$ --- свободная нечётная супералгебра Ли, порождённая $X$.

\begin{lemma}\label{lem:M}
Множество $M$ является базисом $\LL$.
\end{lemma}

\begin{proof}
По замечанию~\ref{rem:parity} $\LL$ --- свободная супералгебра Ли, порождённая множеством $X$, в котором чётность элемента $x\in X$ равна $\rp x$.
Хорошо известно, что свободная супералгебра Ли над полем характеристики нуль, порождённая линейно упорядоченным множеством, имеет базис, состоящий из элементов множества Холла (в частности, из определённых выше базисных коммутаторов) и квадратов $\{v,v\}$ его нечётных элементов $v$, см.~\cite{S86,BMPZ,BKLM} (ср.~\cite[раздел~3]{K17}).
% NOTE FOR THE AUTHORS: [S86, BMPZ, BKLM] usually prove this for specific orderings (regular /
% Lyndon--Shirshov words); please check that the cited statements cover basic commutators of Hall
% with the order used here, or add a sentence explaining why the result holds for every Hall set.
Элемент $v\in\LL$ нечётен относительно обращённой чётности тогда и только тогда, когда $|v|=0$; значит, этот базис есть в точности $M$.
\end{proof}

\subsection*{Множество \texorpdfstring{$U$}{U}}
Пусть
$$
U=\bigl\{\,e_{i_1}^{k_1}e_{i_2}^{k_2}\cdots e_{i_r}^{k_r}\ \bigm|\ r\ge0,\ 1<i_1<i_2<\dots<i_r,\ k_j\ge1,\ \text{и } k_j=1 \text{ при } |e_{i_j}|=1\,\bigr\}
$$
--- множество суперкоммутативных одночленов от элементов $M\setminus\{1\}$ (при $r=0$ пустое произведение равно $1$).
Для элемента $u=e_{i_1}^{k_1}\cdots e_{i_r}^{k_r}$ множества $U$ назовём $\deg_e(u)=\sum_j k_j$ его \emph{$e$-степенью}; элементы $M$ --- это элементы $U$ $e$-степени $1$ вместе с $1$.
Элементы $M$ и $U$ мы рассматриваем как элементы $\G$ (вычисляя слова в $\G$).
Поскольку умножение в $\G$ суперкоммутативно и $ee=0$ для нечётного $e$, произведение двух элементов $U$ с точностью до знака снова является элементом $U$ или нулём.

Основная цель статьи --- доказать, что $U$ является базисом $\G$.
Начнём с более простой половины.

\begin{lemma}\label{vsp}
Множество $U$ порождает $\G$ как векторное пространство.
\end{lemma}

\begin{proof}
Поскольку $(\G,\{\,,\})$ удовлетворяет \eqref{id:skew} и \eqref{id:jacobi}, тождественное отображение на $X$ продолжается до гомоморфизма нечётных супералгебр Ли $\psi\colon\LL\to\G$; он переводит каждый элемент $M$ в тот же элемент $\G$.
По лемме~\ref{lem:M} $\psi(\LL)$ натянуто на $M\subset U$.

Сначала покажем, что $\{u,v\}\in\spn U$ для всех $u,v\in U$, индукцией по $N=\deg_e(u)+\deg_e(v)$, где положено $\deg_e(1)=0$.
Если $u,v\in M$, то $\{u,v\}=\psi(\{u,v\})\in\psi(\LL)\subset\spn U$.
Пусть $\deg_e(v)\ge2$, так что $v=v_1v_2$, где $v_1,v_2\in U\setminus\{1\}$ и $\deg_e(v_1),\deg_e(v_2)<\deg_e(v)$.
По \eqref{id:leibniz}
$$
\{u,v_1v_2\}=\{u,v_1\}v_2+(-1)^{\rp u|v_1|}v_1\{u,v_2\}+(-1)^{\rp u}\{1,u\}v_1v_2 .
$$
Для каждой из пар $(u,v_1)$, $(u,v_2)$ и $(1,u)$ сумма $e$-степеней меньше $N$ (для последней пары она равна $\deg_e(u)<\deg_e(u)+\deg_e(v)$), поэтому $\{u,v_1\}$, $\{u,v_2\}$ и $\{1,u\}$ лежат в $\spn U$ по предположению индукции.
Так как $\spn U$ замкнуто относительно умножения, $\{u,v\}\in\spn U$.
Наконец, если $\deg_e(v)\le1$ и $\deg_e(u)\ge2$, то $\{u,v\}=-(-1)^{\rp u\rp v}\{v,u\}$ по \eqref{id:skew}, а $\{v,u\}\in\spn U$ по предыдущему случаю, применённому к паре $(v,u)$ с тем же значением $N$.

Итак, $\spn U$ замкнуто относительно скобки, а поскольку $U$ замкнуто относительно умножения с точностью до знака, --- и относительно умножения.
Так как $X\subset M\subset U$, подпространство $\spn U$ является подалгеброй $\G$, содержащей образующие; следовательно, $\spn U=\G$.

Заметим, что число букв в слове не годится здесь в качестве параметра индукции, поскольку оно не сохраняется определяющими соотношениями $\G$: соотношение $1a=a$ его уменьшает, а последний член \eqref{id:leibniz} увеличивает.
Например, разложение четырёхбуквенного слова $\{x_1,x_2x_3\}$ содержит слагаемое $(-1)^{\rp{x_1}}\{1,x_1\}x_2x_3$ --- ненулевое кратное пятибуквенного элемента $\{x_1,1\}x_2x_3\in U$.
\end{proof}

\section{Модельная алгебра}\label{sec:model}

В этом разделе мы строим структуру обобщённой алгебры Герстенхабера на свободной суперкоммутативной супералгебре, порождённой $M\setminus\{1\}$.

Пусть $\LL'=\spn(M\setminus\{1\})\subset\LL$, и пусть $\Sm$ --- свободная суперкоммутативная супералгебра, порождённая суперпространством $\LL'$, т.~е. тензорное произведение симметрической алгебры пространства $\LL'_0$ и внешней алгебры пространства $\LL'_1$.
Одночлены от элементов базиса $M\setminus\{1\}$ пространства $\LL'$, т.~е. элементы множества $U$, образуют базис $\Sm$.
Положим $\Lh=\F1\oplus\LL'\subset\Sm$.
Линейное отображение $\LL\to\Lh$, тождественное на $\LL'$ и переводящее базисный элемент $1\in M$ в единицу $1\in\Sm$, --- изоморфизм суперпространств (по лемме~\ref{lem:M}); с его помощью перенесём скобку $\LL$ на $\Lh$.
Тем самым $\Lh$ --- нечётная супералгебра Ли с базисом $M$, в которой $1$ --- наименьший базисный элемент, и $\Lh$ порождает $\Sm$ как унитальную алгебру.
Для $u\in\Lh$ положим $\D(u)=\{1,u\}\in\Lh$; заметим, что $\D(1)=\{1,1\}$ --- нечётный элемент множества $M\setminus\{1\}$.

\subsection*{Обозначения}
Для произведения $a=a_1a_2\cdots a_m$ однородных элементов $\Sm$ и $1\le i\le m$ будем писать
$$
a_{<i}=a_1\cdots a_{i-1},\qquad a_{>i}=a_{i+1}\cdots a_m,\qquad a_{\hat\imath}=a_{<i}a_{>i}
$$
(пустые произведения равны $1$).
По \eqref{id:comm} $a_{<i}a_i=(-1)^{|a_i||a_{<i}|}a_ia_{<i}$ и $a=(-1)^{|a_i||a_{<i}|}a_ia_{\hat\imath}=(-1)^{|a_i||a_{>i}|}a_{\hat\imath}a_i$.

\subsection*{Скобка}
Для $a=a_1\cdots a_m$ и $b=b_1\cdots b_n$, где $m,n\ge1$ и $a_i,b_j\in\Lh$ однородны, положим
\begin{equation}\label{Ddef}
\D(a)=\sum_{i=1}^m(-1)^{|a_{<i}|}\,a_{<i}\{1,a_i\}a_{>i}-(m-1)\{1,1\}\,a
\end{equation}
и
\begin{multline}\label{brdef}
\{a,b\}^*=\sum_{i=1}^m\sum_{j=1}^n(-1)^{|a_i||a_{>i}|+|b_j||b_{<j}|}\,a_{\hat\imath}\{a_i,b_j\}b_{\hat\jmath}\\
-(m-1)\,a\D(b)+(n-1)(-1)^{\rp a}\D(a)b-(m-1)(n-1)\,a\D(1)b .
\end{multline}
Здесь $\{a_i,b_j\}$ --- скобка в $\Lh$, так что $\{a,b\}^*$ выражена через скобку $\LL$ и умножение $\Sm$.
При $m=1$ и $a_1=1$ формула \eqref{Ddef} даёт $\D(1)=\{1,1\}$, в согласии с введённым ранее обозначением, а формула \eqref{brdef} при $m=1$, $a_1=1$ даёт $\{1,b\}^*=\D(b)$ (см. лемму~\ref{lem:welldef}).
Первая сумма в \eqref{brdef} --- это скобка, получаемая переносом $a_i$ в правый конец $a$ и $b_j$ в левый конец $b$ с последующим взятием скобки двух соседних множителей; остальные члены --- поправки, содержащие $\D$.

\begin{lemma}\label{lem:welldef}
Правые части формул \eqref{Ddef} и \eqref{brdef} не зависят от выбора разложений $a=a_1\cdots a_m$ и $b=b_1\cdots b_n$ на множители $a_i,b_j\in\Lh$.
Следовательно, \eqref{Ddef} задаёт линейное отображение $\D\colon\Sm\to\Sm$, а \eqref{brdef} --- билинейное отображение $\{\,,\}^*\colon\Sm\times\Sm\to\Sm$.
Кроме того,
\begin{enumerate}
\item[(i)] $\{u,v\}^*=\{u,v\}$ и $\D(u)=\{1,u\}$ для $u,v\in\Lh$;
\item[(ii)] $\{1,a\}^*=\D(a)$ для всех $a\in\Sm$;
\item[(iii)] $\D(bc)=\D(b)c+(-1)^{|b|}b\D(c)-\D(1)bc$ для всех $b,c\in\Sm$;
\item[(iv)] $|\{a,b\}^*|=|a|+|b|+1$ и $|\D(a)|=\rp a$ для всех однородных $a,b\in\Sm$.
\end{enumerate}
\end{lemma}

\begin{proof}
Зафиксируем $m$ и $n$.
Правая часть \eqref{Ddef} --- полилинейная функция $R_m(a_1,\dots,a_m)$ от $a_1,\dots,a_m\in\Lh$, а правая часть \eqref{brdef} --- полилинейная функция $R_{m,n}(a_1,\dots,a_m;b_1,\dots,b_n)$.
Мы утверждаем, что эти функции
\begin{enumerate}
\item[(a)] суперсимметричны: перестановка двух соседних аргументов $a_k,a_{k+1}$ (соответственно $b_k,b_{k+1}$) умножает значение на $(-1)^{|a_k||a_{k+1}|}$ (соответственно на $(-1)^{|b_k||b_{k+1}|}$);
\item[(b)] согласованы с удалением единиц: $R_m(a_1,\dots,a_{m-1},1)=R_{m-1}(a_1,\dots,a_{m-1})$, $R_{m,n}(a_1,\dots,a_{m-1},1;b)=R_{m-1,n}(a_1,\dots,a_{m-1};b)$ при $m\ge2$, и аналогично для второй группы аргументов.
\end{enumerate}
Проверим (a) для $R_{m,n}$ и аргументов $a_k,a_{k+1}$; остальные случаи аналогичны или проще.
Произведения $a$, $a\D(b)$ и $a\D(1)b$ умножаются на $(-1)^{|a_k||a_{k+1}|}$, и то же верно для $\D(a)b$, как только (a) известно для $R_m$.
В двойной сумме члены с $i\ne k,k+1$ умножаются на $(-1)^{|a_k||a_{k+1}|}$, поскольку так меняется $a_{\hat\imath}$, а показатель $|a_i||a_{>i}|$ не меняется.
Член с $i=k$ исходной суммы равен $(-1)^{|a_k|(|a_{k+1}|+|a_{>k+1}|)}a_{<k}a_{k+1}a_{>k+1}\{a_k,b_j\}b_{\hat\jmath}$, а член с $i=k+1$ суммы для переставленных аргументов равен $(-1)^{|a_k||a_{>k+1}|}a_{<k}a_{k+1}a_{>k+1}\{a_k,b_j\}b_{\hat\jmath}$; они отличаются множителем $(-1)^{|a_k||a_{k+1}|}$, и то же верно для пары членов с $i=k+1$ и $i=k$.
Для $R_m$ члены с $i\ne k,k+1$ и последний член \eqref{Ddef} разбираются так же (чётность $a_{<i}$ не меняется).
Член с $i=k$ исходной суммы равен $(-1)^{|a_{<k}|}a_{<k}\{1,a_k\}a_{k+1}a_{>k+1}$, а член с $i=k+1$ суммы для переставленных аргументов равен
$$
(-1)^{|a_{<k}|+|a_{k+1}|}a_{<k}a_{k+1}\{1,a_k\}a_{>k+1}=(-1)^{|a_{<k}|+|a_{k+1}|+|a_{k+1}|(|a_k|+1)}a_{<k}\{1,a_k\}a_{k+1}a_{>k+1},
$$
так как $|\{1,a_k\}|=|a_k|+1$; здесь знак, происходящий от изменения множителя $(-1)^{|a_{<i}|}$, компенсируется знаком, возникающим при переносе $\{1,a_k\}$ через $a_{k+1}$, и два члена отличаются множителем $(-1)^{|a_k||a_{k+1}|}$.
Пара членов с $i=k+1$ и $i=k$ разбирается так же.

Для проверки (b) пусть $a_m=1$ и $a'=a_1\cdots a_{m-1}$.
В \eqref{Ddef} член с $i=m$ равен $(-1)^{|a'|}a'\{1,1\}=\{1,1\}a'$ (напомним, что $\{1,1\}$ нечётен), так что $R_m(a_1,\dots,a_{m-1},1)=R_{m-1}(a_1,\dots,a_{m-1})+\{1,1\}a'-\{1,1\}a'$.
В \eqref{brdef} члены с $i=m$ дают вклад $a'\sum_{j}(-1)^{|b_j||b_{<j}|}\{1,b_j\}b_{\hat\jmath}$, а поскольку $(-1)^{|b_j||b_{<j}|}\{1,b_j\}b_{\hat\jmath}=(-1)^{|b_{<j}|}b_{<j}\{1,b_j\}b_{>j}$, этот вклад равен $a'\bigl(\D(b)+(n-1)\D(1)b\bigr)$ по \eqref{Ddef}; члены с $i<m$ дают двойную сумму для $a'$.
Остальные члены \eqref{brdef} для $a$ и для $a'$ отличаются на $-a'\D(b)-(n-1)a'\D(1)b$, что сокращается с дополнительным вкладом двойной суммы.
Проверка для $b_n=1$ аналогична: положив $b'=b_1\cdots b_{n-1}$ и используя $\{a_i,1\}=-(-1)^{\rp{a_i}}\{1,a_i\}$ (тождество \eqref{id:skew} в $\Lh$), $a_{<i}\{1,a_i\}a_{>i}=(-1)^{\rp{a_i}|a_{>i}|}a_{\hat\imath}\{1,a_i\}$ и \eqref{Ddef}, находим, что члены с $j=n$ двойной суммы дают вклад $(-1)^{|a|}\bigl(\D(a)+(m-1)\D(1)a\bigr)b'$, который сокращается с разностью $(-1)^{\rp a}\D(a)b'-(m-1)\,a\D(1)b'$ между остальными членами \eqref{brdef} для $b$ и для $b'$.

Теперь определим $\D$ на базисе $U$ пространства $\Sm$ и $\{\,,\}^*$ на $U\times U$ формулами \eqref{Ddef} и \eqref{brdef}, используя для $u=e_{i_1}^{k_1}\cdots e_{i_r}^{k_r}\in U$ разложение на $\deg_e(u)$ множителей из $M\setminus\{1\}$, взятых в неубывающем порядке, а для $u=1$ --- разложение с $m=1$, $a_1=1$; продолжим по линейности.
Тогда \eqref{Ddef} и \eqref{brdef} выполняются для произвольных разложений: обе части полилинейны по множителям, поэтому достаточно рассмотреть множители из базиса $M$ пространства $\Lh$; по (a) множители можно переставить в неубывающем порядке, причём обе части меняются на один и тот же знак; если некоторый нечётный базисный элемент встречается дважды, левая часть равна нулю, и правая часть тоже, по (a), поскольку она меняет знак при перестановке двух равных нечётных аргументов; по (b) множители, равные $1$, можно удалить, и остаётся разложение, использованное в определении.

Утверждения (i), (ii) и (iv) следуют непосредственно из \eqref{Ddef} и \eqref{brdef} при $m=n=1$, соответственно при $m=1$, $a_1=1$ (с использованием установленного выше тождества $(-1)^{|b_j||b_{<j}|}\{1,b_j\}b_{\hat\jmath}=(-1)^{|b_{<j}|}b_{<j}\{1,b_j\}b_{>j}$ и $(-1)^{\rp 1}=-1$).
Для (iii) применим \eqref{Ddef} к разложению $bc=b_1\cdots b_nc_1\cdots c_p$: члены, отвечающие множителям $b$, дают $(\D(b)+(n-1)\D(1)b)c$, члены, отвечающие множителям $c$, дают $(-1)^{|b|}b(\D(c)+(p-1)\D(1)c)=(-1)^{|b|}b\D(c)+(p-1)\D(1)bc$, а последний член равен $-(n+p-1)\D(1)bc$.
\end{proof}

Проверим теперь тождества \eqref{id:skew} и \eqref{id:leibniz} для $\{\,,\}^*$.
В доказательствах обозначим через $\{a,b\}_0$ двойную сумму в \eqref{brdef}, так что
\begin{equation}\label{brsplit}
\{a,b\}^*=\{a,b\}_0-(m-1)\,a\D(b)+(n-1)(-1)^{\rp a}\D(a)b-(m-1)(n-1)\,a\D(1)b .
\end{equation}

\begin{lemma}\label{lem:skew}
Для всех $a,b\in\Sm$ выполнено $\{a,b\}^*=-(-1)^{\rp a\rp b}\{b,a\}^*$.
\end{lemma}

\begin{proof}
Пусть $a=a_1\cdots a_m$, $b=b_1\cdots b_n$, где $a_i,b_j\in\Lh$.
Рассмотрим член $\{b,a\}_0$, отвечающий паре $(j,i)$:
$$
(-1)^{|b_j||b_{>j}|+|a_i||a_{<i}|}\,b_{\hat\jmath}\{b_j,a_i\}a_{\hat\imath}.
$$
По \eqref{id:skew} в $\Lh$ имеем $\{b_j,a_i\}=-(-1)^{\rp{a_i}\rp{b_j}}\{a_i,b_j\}$, а по \eqref{id:comm}
$$
b_{\hat\jmath}\{a_i,b_j\}a_{\hat\imath}=(-1)^{|a_{\hat\imath}|(|b_{\hat\jmath}|+\rp{a_i}+|b_j|)+|b_{\hat\jmath}|(\rp{a_i}+|b_j|)}\,a_{\hat\imath}\{a_i,b_j\}b_{\hat\jmath},
$$
поскольку $|\{a_i,b_j\}|=\rp{a_i}+|b_j|$.
Используя $\rp a=|a_{\hat\imath}|+\rp{a_i}$, $\rp b=|b_{\hat\jmath}|+\rp{b_j}$, $|b_{\hat\jmath}|=|b|+|b_j|$, $|a_{\hat\imath}|=|a|+|a_i|$ и $|a_i|+\rp{a_i}=|b_j|+\rp{b_j}=1$, прямым сравнением показателей получаем, что этот член равен
$$
-(-1)^{\rp a\rp b}\,(-1)^{|a_i||a_{>i}|+|b_j||b_{<j}|}\,a_{\hat\imath}\{a_i,b_j\}b_{\hat\jmath},
$$
т.~е. $-(-1)^{\rp a\rp b}$, умноженному на $(i,j)$-член $\{a,b\}_0$.
Следовательно, $\{b,a\}_0=-(-1)^{\rp a\rp b}\{a,b\}_0$.

Для остальных членов \eqref{brsplit} воспользуемся $|\D(a)|=\rp a$, $|\D(1)|=1$ и \eqref{id:comm}:
$$
b\D(a)=(-1)^{|b|\rp a}\D(a)b,\qquad \D(b)a=(-1)^{\rp b|a|}a\D(b),\qquad b\D(1)a=-(-1)^{\rp a\rp b}a\D(1)b .
$$
Поэтому
\begin{multline*}
-(n-1)\,b\D(a)+(m-1)(-1)^{\rp b}\D(b)a-(n-1)(m-1)\,b\D(1)a\\
=-(-1)^{\rp a\rp b}\Bigl(-(m-1)\,a\D(b)+(n-1)(-1)^{\rp a}\D(a)b-(m-1)(n-1)\,a\D(1)b\Bigr),
\end{multline*}
поскольку $(-1)^{\rp b(|a|+1)}=(-1)^{\rp a\rp b}$ и $(-1)^{|b|\rp a}=(-1)^{\rp a\rp b+\rp a}$.
Складывая оба вычисления, получаем утверждение леммы.
\end{proof}

\begin{lemma}\label{lem:leibniz}
Для всех $a,b,c\in\Sm$ выполнено
$$
\{a,bc\}^*=\{a,b\}^*c+(-1)^{\rp a|b|}b\{a,c\}^*+(-1)^{\rp a}\D(a)bc .
$$
\end{lemma}

\begin{proof}
Пусть $a=a_1\cdots a_m$, $b=b_1\cdots b_n$, $c=c_1\cdots c_p$, где $a_i,b_j,c_k\in\Lh$, и применим \eqref{brdef} к разложению $bc=b_1\cdots b_nc_1\cdots c_p$.
Члены $\{a,bc\}_0$, отвечающие множителям $b$, дают $\{a,b\}_0c$.
Член, отвечающий множителю $c_k$, равен
$$
(-1)^{|a_i||a_{>i}|+|c_k|(|b|+|c_{<k}|)}\,a_{\hat\imath}\{a_i,c_k\}\,b\,c_{\hat k}
=(-1)^{\rp a|b|}\,(-1)^{|a_i||a_{>i}|+|c_k||c_{<k}|}\,b\,a_{\hat\imath}\{a_i,c_k\}c_{\hat k},
$$
поскольку $|a_{\hat\imath}\{a_i,c_k\}|=\rp a+|c_k|$; значит, эти члены дают $(-1)^{\rp a|b|}b\{a,c\}_0$.
Итак,
$$
\{a,bc\}_0=\{a,b\}_0c+(-1)^{\rp a|b|}b\{a,c\}_0 .
$$
Остаётся сравнить члены, содержащие $\D$.
По \eqref{brsplit} в левой части они равны
$$
-(m-1)\,a\D(bc)+(n+p-1)(-1)^{\rp a}\D(a)bc-(m-1)(n+p-1)\,a\D(1)bc,
$$
а в правой части ---
\begin{multline*}
\bigl(-(m-1)\,a\D(b)+(n-1)(-1)^{\rp a}\D(a)b-(m-1)(n-1)\,a\D(1)b\bigr)c\\
+(-1)^{\rp a|b|}b\bigl(-(m-1)\,a\D(c)+(p-1)(-1)^{\rp a}\D(a)c-(m-1)(p-1)\,a\D(1)c\bigr)+(-1)^{\rp a}\D(a)bc .
\end{multline*}
Используя \eqref{id:comm} вместе с $|\D(a)|=\rp a$ и $|\D(1)|=1$, имеем
$$
(-1)^{\rp a|b|}ba\D(c)=(-1)^{|b|}ab\D(c),\quad (-1)^{\rp a|b|}b\D(a)c=\D(a)bc,\quad (-1)^{\rp a|b|}ba\D(1)c=a\D(1)bc,
$$
так что правая часть равна
$$
-(m-1)\,a\bigl(\D(b)c+(-1)^{|b|}b\D(c)\bigr)+(n+p-1)(-1)^{\rp a}\D(a)bc-(m-1)(n+p-2)\,a\D(1)bc,
$$
что совпадает с левой частью по лемме~\ref{lem:welldef}(iii).
\end{proof}

Для доказательства тождества Якоби воспользуемся следующей общей формулой для якобиатора.

\begin{lemma}\label{jacobiator}
Пусть $(V,\cdot,\{\,,\})$ --- супералгебра с единицей, удовлетворяющая \eqref{id:deg}, \eqref{id:comm}, \eqref{id:skew} и \eqref{id:leibniz}, и пусть
$$
J(a,b,c)=\{a,\{b,c\}\}-\{\{a,b\},c\}-(-1)^{\rp a\rp b}\{b,\{a,c\}\}
$$
для $a,b,c\in V$.
Тогда
\begin{enumerate}
\item[(i)] $J(a,b,c)=-(-1)^{\rp a\rp b}J(b,a,c)=-(-1)^{\rp b\rp c}J(a,c,b)$;
\item[(ii)] $J(a,b,xy)=J(a,b,x)\,y+(-1)^{(|a|+|b|)|x|}\,x\,J(a,b,y)-J(a,b,1)\,xy$ для всех $a,b,x,y\in V$.
\end{enumerate}
\end{lemma}

\begin{proof}
(i) По \eqref{id:skew} $\{\{b,a\},c\}=-(-1)^{\rp a\rp b}\{\{a,b\},c\}$, откуда
$$
-(-1)^{\rp a\rp b}J(b,a,c)=-(-1)^{\rp a\rp b}\{b,\{a,c\}\}-\{\{a,b\},c\}+\{a,\{b,c\}\}=J(a,b,c).
$$
Аналогично, $\{a,\{c,b\}\}=-(-1)^{\rp b\rp c}\{a,\{b,c\}\}$, $\{\{a,c\},b\}=-(-1)^{\rp b(\rp a+\rp c)}\{b,\{a,c\}\}$ и $\{c,\{a,b\}\}=-(-1)^{\rp c(\rp a+\rp b)}\{\{a,b\},c\}$, поскольку $\rp{\{a,c\}}=\rp a+\rp c$ и $\rp{\{a,b\}}=\rp a+\rp b$; подставляя это в $J(a,c,b)$, получаем $J(a,c,b)=-(-1)^{\rp b\rp c}J(a,b,c)$.

(ii) Раскроем три члена $J(a,b,xy)$ по \eqref{id:leibniz}, используя $|\{b,x\}|=\rp b+|x|$, $|\D(b)|=\rp b$, $\rp{\{a,b\}}=\rp a+\rp b$ и $(-1)^{(\rp a+\rp b)|x|}=(-1)^{(|a|+|b|)|x|}$:
\begin{align*}
\{a,\{b,xy\}\}={}&\{a,\{b,x\}\}y+(-1)^{\rp a(|x|+\rp b)}\{b,x\}\{a,y\}+(-1)^{\rp a}\D(a)\{b,x\}y\\
&+(-1)^{\rp b|x|}\Bigl(\{a,x\}\{b,y\}+(-1)^{\rp a|x|}x\{a,\{b,y\}\}+(-1)^{\rp a}\D(a)x\{b,y\}\Bigr)\\
&+(-1)^{\rp b}\Bigl(\{a,\D(b)\}xy+(-1)^{\rp a\rp b}\D(b)\{a,xy\}+(-1)^{\rp a}\D(a)\D(b)xy\Bigr),\\
\{\{a,b\},xy\}={}&\{\{a,b\},x\}y+(-1)^{(|a|+|b|)|x|}x\{\{a,b\},y\}+(-1)^{\rp a+\rp b}\D(\{a,b\})xy,\\
\{b,\{a,xy\}\}={}&\{b,\{a,x\}\}y+(-1)^{\rp b(|x|+\rp a)}\{a,x\}\{b,y\}+(-1)^{\rp b}\D(b)\{a,x\}y\\
&+(-1)^{\rp a|x|}\Bigl(\{b,x\}\{a,y\}+(-1)^{\rp b|x|}x\{b,\{a,y\}\}+(-1)^{\rp b}\D(b)x\{a,y\}\Bigr)\\
&+(-1)^{\rp a}\Bigl(\{b,\D(a)\}xy+(-1)^{\rp a\rp b}\D(a)\{b,xy\}+(-1)^{\rp b}\D(b)\D(a)xy\Bigr),
\end{align*}
где, кроме того, $\{a,xy\}$ и $\{b,xy\}$ следует раскрыть по \eqref{id:leibniz}.
Подставим эти разложения в $J(a,b,xy)=\{a,\{b,xy\}\}-\{\{a,b\},xy\}-(-1)^{\rp a\rp b}\{b,\{a,xy\}\}$.
Члены $\{a,\{b,x\}\}y$, $-\{\{a,b\},x\}y$, $-(-1)^{\rp a\rp b}\{b,\{a,x\}\}y$ в сумме дают $J(a,b,x)y$, а члены, содержащие $x\{a,\{b,y\}\}$, $x\{\{a,b\},y\}$, $x\{b,\{a,y\}\}$, в сумме дают $(-1)^{(|a|+|b|)|x|}xJ(a,b,y)$.
По \eqref{id:skew} имеем $\{b,1\}=-(-1)^{\rp b}\D(b)$, $\{a,1\}=-(-1)^{\rp a}\D(a)$ и $\{\{a,b\},1\}=-(-1)^{\rp a+\rp b}\D(\{a,b\})$, так что
$$
J(a,b,1)=-(-1)^{\rp b}\{a,\D(b)\}+(-1)^{\rp a+\rp b}\D(\{a,b\})+(-1)^{\rp a\rp b+\rp a}\{b,\D(a)\};
$$
поэтому члены $(-1)^{\rp b}\{a,\D(b)\}xy$, $-(-1)^{\rp a+\rp b}\D(\{a,b\})xy$, $-(-1)^{\rp a\rp b+\rp a}\{b,\D(a)\}xy$ в сумме дают $-J(a,b,1)xy$.
Оставшиеся шестнадцать членов сокращаются попарно: после раскрытия $\{a,xy\}$ и $\{b,xy\}$ их сумма равна
\begin{gather*}
\bigl((-1)^{\rp a(|x|+\rp b)}-(-1)^{\rp a\rp b+\rp a|x|}\bigr)\{b,x\}\{a,y\}
+\bigl((-1)^{\rp b|x|}-(-1)^{\rp a\rp b+\rp b(|x|+\rp a)}\bigr)\{a,x\}\{b,y\}\\
+\bigl((-1)^{\rp a}-(-1)^{\rp a}\bigr)\D(a)\{b,x\}y
+\bigl((-1)^{\rp b|x|+\rp a}-(-1)^{\rp a+\rp b|x|}\bigr)\D(a)x\{b,y\}\\
+\bigl((-1)^{\rp b+\rp a\rp b}-(-1)^{\rp a\rp b+\rp b}\bigr)\D(b)\{a,x\}y\\
+\bigl((-1)^{\rp b+\rp a\rp b+\rp a|x|}-(-1)^{\rp a\rp b+\rp a|x|+\rp b}\bigr)\D(b)x\{a,y\}\\
+\bigl((-1)^{\rp b+\rp a}-(-1)^{\rp a+\rp b}\bigr)\D(a)\D(b)xy
+\bigl((-1)^{\rp b+\rp a\rp b+\rp a}-(-1)^{\rp a\rp b+\rp a+\rp b}\bigr)\D(b)\D(a)xy,
\end{gather*}
и каждый из восьми коэффициентов равен нулю.
\end{proof}

\begin{proposition}\label{prop:jacobi}
Скобка $\{\,,\}^*$ удовлетворяет нечётному тождеству Якоби \eqref{id:jacobi} на $\Sm$.
\end{proposition}

\begin{proof}
По леммам~\ref{lem:welldef}, \ref{lem:skew} и \ref{lem:leibniz} супералгебра $(\Sm,\cdot,\{\,,\}^*)$ удовлетворяет условиям леммы~\ref{jacobiator}; пусть $J$ --- её якобиатор.
Нужно показать, что $J(a,b,c)=0$ для всех $a,b,c\in\Sm$.

Сначала заметим, что при фиксированных $a,b\in\Sm$ из $J(a,b,u)=0$ для всех $u\in\Lh$ следует $J(a,b,c)=0$ для всех $c\in\Sm$.
Действительно, $\Sm$ натянуто на произведения $c=c_1\cdots c_m$ элементов $\Lh$, и по лемме~\ref{jacobiator}(ii)
\begin{multline*}
J(a,b,c_1\cdots c_m)=J(a,b,c_1\cdots c_{m-1})\,c_m\\
+(-1)^{(|a|+|b|)|c_1\cdots c_{m-1}|}\,c_1\cdots c_{m-1}\,J(a,b,c_m)-J(a,b,1)\,c_1\cdots c_m,
\end{multline*}
так что утверждение следует индукцией по $m$, поскольку $1,c_m\in\Lh$.

Пусть теперь $u,v,w\in\Lh$.
По лемме~\ref{lem:welldef}(i) $J(u,v,w)$ --- якобиатор нечётной супералгебры Ли $\Lh$, поэтому $J(u,v,w)=0$.
По предыдущему абзацу $J(u,v,c)=0$ для всех $c\in\Sm$.
По лемме~\ref{jacobiator}(i) $J(u,c,v)=-(-1)^{\rp c\rp v}J(u,v,c)=0$ для $u,v\in\Lh$, $c\in\Sm$; применяя предыдущий абзац ещё раз, получаем $J(u,b,c)=0$ для всех $u\in\Lh$ и $b,c\in\Sm$.
Наконец, по лемме~\ref{jacobiator}(i) $J(a,b,u)=-(-1)^{\rp b\rp u}J(a,u,b)=(-1)^{\rp b\rp u+\rp a\rp u}J(u,a,b)=0$ для $a,b\in\Sm$, $u\in\Lh$, и первый абзац даёт $J(a,b,c)=0$ для всех $a,b,c\in\Sm$.
\end{proof}

Подводя итог, получаем

\begin{theorem}\label{thm:model}
$(\Sm,\cdot,\{\,,\}^*)$ --- унитальная обобщённая алгебра Герстенхабера, ограничение скобки которой на $\Lh$ есть скобка свободной нечётной супералгебры Ли $\LL$, а отображение $\D$ задаётся формулой \eqref{Ddef}.
Она порождена множеством $X$ как унитальная обобщённая алгебра Герстенхабера.
\end{theorem}

\begin{proof}
Тождества \eqref{id:deg} и \eqref{id:comm} выполнены в $\Sm$ по построению и по лемме~\ref{lem:welldef}(iv); \eqref{id:skew}, \eqref{id:leibniz} и \eqref{id:jacobi} --- это лемма~\ref{lem:skew}, лемма~\ref{lem:leibniz} (заметим, что $\D(a)=\{1,a\}^*$ по лемме~\ref{lem:welldef}(ii)) и предложение~\ref{prop:jacobi}.
Подалгебра $\Sm$, порождённая $X$, содержит $\Lh$, поскольку $\LL$ порождена $X$ как нечётная супералгебра Ли, а $\{\,,\}^*$ совпадает на $\Lh$ со скобкой $\LL$; а $\Lh$ порождает $\Sm$ относительно умножения.
\end{proof}

\section{Основные результаты}

\begin{theorem}\label{basis}
Множество $U$ является базисом свободной унитальной обобщённой алгебры Герстенхабера $\G$, порождённой $X$.
Более того, гомоморфизм $\varphi\colon\G\to\Sm$, тождественный на $X$, является изоморфизмом обобщённых алгебр Герстенхабера.
\end{theorem}

\begin{proof}
По теореме~\ref{thm:model} и универсальному свойству $\G$ тождественное отображение на $X$ продолжается до гомоморфизма $\varphi\colon\G\to\Sm$ унитальных обобщённых алгебр Герстенхабера, сюръективного, поскольку $\Sm$ порождена $X$.
Каждый элемент $e\in M$ --- слово от образующих относительно скобки, а $\varphi$ согласован со скобками; поэтому $\varphi(e)=e$ для $e\in M$ и, следовательно, $\varphi(u)=u$ для всех $u\in U$.
Так как $U$ --- базис $\Sm$, множество $U\subset\G$ линейно независимо, а по лемме~\ref{vsp} оно порождает $\G$.
Итак, $U$ --- базис $\G$, а $\varphi$, будучи сюръекцией, биективно переводящей базис в базис, --- изоморфизм.
\end{proof}

\begin{corollary}\label{cor:11}
Элемент $\D(1)=\{1,1\}$ отличен от нуля в $\G$, и $\D$ не является дифференцированием алгебры $(\G,\cdot)$.
Точнее, элементы $\{1,1\}$, $\{1,1\}x_i$ и $\{1,x_i\}$, $i\ge1$, линейно независимы в $\G$.
\end{corollary}

\begin{proof}
Слова $\{1,1\}$ (квадрат чётного хорошего слова $1$), $\{x_i,1\}$ (хорошее слово) и одночлены $\{1,1\}x_i$ принадлежат $U$; заметим, что $\{1,x_i\}=-(-1)^{\rp{x_i}}\{x_i,1\}$.
Второе утверждение следует из замечания~\ref{rem:D}(a).
Первое утверждение можно увидеть и непосредственно, без теоремы~\ref{basis}: по универсальному свойству $\G$ существует гомоморфизм из $\G$ в алгебру $V$ примера~\ref{ex:dim2}, переводящий $1$ в $1$ и каждый $x_i$ в $0$, и он переводит $\{1,1\}$ в $\varepsilon\ne0$.
\end{proof}

Следуя~\cite{K17}, для $n\ge1$ рассмотрим свободную унитальную обобщённую алгебру Герстенхабера $\G$, порождённую $X=\{1,x_1,\dots,x_n\}$ (с произвольными чётностями $x_1,\dots,x_n$), и обозначим через $\mathrm{GenGer}(n)$ подпространство $\G$, натянутое на те элементы базиса $U$, которые полилинейны по $1,x_1,\dots,x_n$, т.~е. в которых каждая из букв $1,x_1,\dots,x_n$ встречается ровно один раз.
Заметим, что $\G$ градуирована мультистепенью по $x_1,\dots,x_n$, но не числом вхождений буквы $1$: соотношение $1a=a$ и последний член \eqref{id:leibniz} не однородны по $1$.
Поэтому $\mathrm{GenGer}(n)$ определяется как оболочка подмножества базиса $U$ (эквивалентно, через изоморфизм $\varphi$ теоремы~\ref{basis}, --- одночленов $\Sm$), а не как однородная компонента $\G$; это подпространство компоненты $\G$ мультистепени $(1,\dots,1)$ по $x_1,\dots,x_n$.

\begin{theorem}\label{dimen}
$\dim\mathrm{GenGer}(n)=n\cdot n!$.
\end{theorem}

\begin{proof}
Элемент $U$ --- это произведение $e_{i_1}\cdots e_{i_r}$ элементов $M\setminus\{1\}$, и он полилинеен по $1,x_1,\dots,x_n$ тогда и только тогда, когда его множители --- хорошие слова (квадраты $\{v,v\}$ содержат каждую букву $v$ дважды), множества букв которых попарно не пересекаются, полилинейны и в объединении дают $X$.
Таким образом, полилинейные элементы $U$ находятся во взаимно однозначном соответствии с парами, состоящими из разбиения $\pi$ множества $X$ на блоки, ни один из которых не равен $\{1\}$ (множитель $1$ в $U$ исключён), и семейства полилинейных хороших слов $w_B$, $B\in\pi$, где множество букв слова $w_B$ равно $B$.

Полилинейные хорошие слова с данным множеством $B$ из $k$ букв образуют, по лемме~\ref{lem:M}, базис полилинейной компоненты свободной нечётной супералгебры Ли, порождённой $B$; поскольку хорошие слова определены чисто комбинаторно, их число не зависит от чётностей букв и равно размерности $(k-1)!$ полилинейной компоненты свободной алгебры Ли ранга $k$ (см.~\cite{Reu93}).
Следовательно,
$$
\dim\mathrm{GenGer}(n)=\sum_{\pi}\ \prod_{B\in\pi}(|B|-1)!,
$$
где $\pi$ пробегает разбиения $X$ без блока $\{1\}$.
Для всякого конечного множества $Y$ имеем $\sum_{\pi}\prod_{B\in\pi}(|B|-1)!=|Y|!$, где сумма берётся по всем разбиениям $Y$: разбиение $Y$ вместе с циклическим порядком на каждом блоке --- это то же самое, что перестановка $Y$ (блоки --- её циклы).
Вычитая разбиения $X$, содержащие блок $\{1\}$, которые соответствуют разбиениям множества $\{x_1,\dots,x_n\}$, получаем $\dim\mathrm{GenGer}(n)=(n+1)!-n!=n\cdot n!$.
\end{proof}

\begin{example}\label{ex:small}
При $n=1$ единственный полилинейный элемент $U$ --- хорошее слово $\{x_1,1\}$, так что $\dim\mathrm{GenGer}(1)=1=1\cdot1!$.
При $n=2$ полилинейные хорошие слова с множеством букв $\{1,x_1,x_2\}$ --- это $\{\{x_1,1\},x_2\}$ и $\{\{x_2,1\},x_1\}$ (слово $\{\{x_2,x_1\},1\}$ не является хорошим, поскольку $x_1>1$), а разбиения $\{1,x_1\}\cup\{x_2\}$ и $\{1,x_2\}\cup\{x_1\}$ дают произведения $\{x_1,1\}x_2$ и $\{x_2,1\}x_1$.
Следовательно, $\mathrm{GenGer}(2)$ имеет базис
$$
\{x_1,1\}x_2,\qquad \{x_2,1\}x_1,\qquad \{\{x_1,1\},x_2\},\qquad \{\{x_2,1\},x_1\},
$$
и $\dim\mathrm{GenGer}(2)=4=3!-2!=2\cdot2!$, в согласии с теоремой~\ref{dimen}.
\end{example}

\begin{remark}
То же рассуждение, применённое к свободной нечётной супералгебре Ли $\LL_0$, порождённой $\{x_1,x_2,\dots\}$, со скобкой, продолженной на $\F1\oplus\LL_0$ нулём ($\{1,\cdot\}=\{\cdot,1\}=0$), даёт $\D=0$ в \eqref{Ddef} и превращает \eqref{brdef} в единственное продолжение скобки $\LL_0$ до бидифференцирования свободной суперкоммутативной супералгебры, порождённой $\LL_0$.
Леммы~\ref{lem:welldef}--\ref{jacobiator} и предложение~\ref{prop:jacobi} остаются верными дословно, а доказательство леммы~\ref{vsp} проходит с $\{1,u\}=0$.
Так восстанавливается хорошо известное описание свободной унитальной алгебры Герстенхабера, порождённой $x_1,x_2,\dots$: её базис образуют элементы $U$, в которых буква $1$ не встречается.
\end{remark}

\begin{thebibliography}{99}

\bibitem{BMPZ}
Bahturin Yu., Mikhalev A., Petrogradsky V., Zaicev M.,
Infinite-dimensional Lie superalgebras,
De Gruyter Expositions in Mathematics, 7. Walter de Gruyter \& Co., Berlin, 1992.

\bibitem{BV81}
Batalin I., Vilkovisky G.,
Gauge algebra and quantization,
Physics Letters B, 102 (1981), 1, 27--31.

\bibitem{p2}
Batalin I., Bering K.,
Odd scalar curvature in anti-Poisson geometry,
Physics Letters B, 663 (2008), 1--2, 132--135.

\bibitem{BLLM}
Bell J., Launois S., Le\'on S\'anchez O., Moosa R.,
Poisson algebras via model theory and differential algebraic geometry,
Journal of the European Mathematical Society, 19 (2017), 7, 2019--2049.

\bibitem{BKLM}
Bokut L., Kang S.-J., Lee K.-H., Malcolmson P.,
Gr\"obner--Shirshov bases for Lie superalgebras and their universal enveloping algebras,
Journal of Algebra, 217 (1999), 2, 461--495.

\bibitem{B06}
Bratchikov A.,
Subalgebras of the Gerstenhaber algebra of differential operators,
Journal of Algebra and its Applications, 16 (2017), 1, 1750005, 6 pp.

\bibitem{p1}
Bruce A., Grabowski J.,
Odd connections on supermanifolds: existence and relation with affine connections,
Journal of Physics A, 53 (2020), 45, 455203, 24 pp.

\bibitem{CK07}
Cantarini N., Kac V.,
Classification of linearly compact simple Jordan and generalized Poisson superalgebras,
Journal of Algebra, 313 (2007), 1, 100--124.

\bibitem{CK10}
Cantarini N., Kac V.,
Classification of linearly compact simple rigid superalgebras,
International Mathematics Research Notices, (2010), 17, 3341--3393.

\bibitem{chi}
Chibrikov E.,
A right normed basis for free Lie algebras and Lyndon--Shirshov words,
Journal of Algebra, 302 (2006), 2, 593--612.

\bibitem{Ger63}
Gerstenhaber M.,
The cohomology structure of an associative ring,
Annals of Mathematics, 78 (1963), 2, 267--288.

\bibitem{GK04}
Ginzburg V., Kaledin D.,
Poisson deformations of symplectic quotient singularities,
Advances in Mathematics, 186 (2004), 1, 1--57.

\bibitem{Hall50}
Hall M.,
A basis for free Lie rings and higher commutators in free groups,
Proceedings of the American Mathematical Society, 1 (1950), 575--581.

\bibitem{Kac77}
Kac V.,
Lie superalgebras,
Advances in Mathematics, 26 (1977), 1, 8--96.

\bibitem{Kan90}
Kantor I.,
Connection between Poisson brackets and Jordan and Lie superalgebras,
Lie theory, differential equations and representation theory (Montreal, 1989), 213--225, Univ. Montr\'eal, Montreal, 1990.

\bibitem{Kan92}
Kantor I.,
Jordan and Lie superalgebras determined by a Poisson algebra,
American Mathematical Society Translations, Series 2, 151 (1992), 55--80.

\bibitem{K10}
Kaygorodov I.,
Generalized Kantor double,
Vestnik Samarskogo Gosudarstvennogo Universiteta, 78 (2010), 4, 42--50.

\bibitem{K17}
Kaygorodov I.,
Algebras of Jordan brackets and generalized Poisson algebras,
Linear and Multilinear Algebra, 65 (2017), 6, 1142--1157.

\bibitem{KSU}
Kaygorodov I., Shestakov I., Umirbaev U.,
Free generic Poisson fields and algebras,
Communications in Algebra, 46 (2018), 4, 1799--1812.

\bibitem{KM92}
King D., McCrimmon K.,
The Kantor construction of Jordan superalgebras,
Communications in Algebra, 20 (1992), 1, 109--126.

\bibitem{Kon03}
Kontsevich M.,
Deformation quantization of Poisson manifolds,
Letters in Mathematical Physics, 66 (2003), 3, 157--216.

\bibitem{K96}
Kosmann-Schwarzbach Y.,
From Poisson algebras to Gerstenhaber algebras,
Annales de l'Institut Fourier (Grenoble), 46 (1996), 5, 1243--1274.

\bibitem{L1}
Li L.-C.,
Classical $r$-matrices and compatible Poisson structures for Lax equations on Poisson algebras,
Communications in Mathematical Physics, 203 (1999), 3, 573--592.

\bibitem{L21}
Lopes S., Solotar A.,
Lie structure on the Hochschild cohomology of a family of subalgebras of the Weyl algebra,
Journal of Noncommutative Geometry, 15 (2021), 4, 1373--1407.

\bibitem{Reu93}
Reutenauer C.,
Free Lie algebras,
London Mathematical Society Monographs. New Series, 7. Oxford University Press, New York, 1993.

\bibitem{She93}
Shestakov I.,
Quantization of Poisson superalgebras and speciality of Jordan Poisson superalgebras,
Algebra and Logic, 32 (1993), 5, 309--317.

\bibitem{shirshov}
Shirshov A. I.,
Selected works of A. I. Shirshov,
Translated from the Russian by M. Bremner and M. V. Kotchetov. Edited by L. A. Bokut, V. Latyshev, I. Shestakov and E. Zelmanov. Contemporary Mathematicians. Birkh\"auser Verlag, Basel, 2009.

\bibitem{S86}
Shtern A.,
Free Lie superalgebras,
Siberian Mathematical Journal, 27 (1986), 1, 136--140.

\bibitem{ber08}
Van den Bergh M.,
Double Poisson algebras,
Transactions of the American Mathematical Society, 360 (2008), 11, 5711--5769.

\bibitem{V16}
Volkov Yu.,
BV differential on Hochschild cohomology of Frobenius algebras,
Journal of Pure and Applied Algebra, 220 (2016), 10, 3384--3402.

\bibitem{V19}
Volkov Yu.,
Gerstenhaber bracket on the Hochschild cohomology via an arbitrary resolution,
Proceedings of the Edinburgh Mathematical Society (2), 62 (2019), 3, 817--836.

\bibitem{W21}
Wang Z.,
Gerstenhaber algebra and Deligne's conjecture on the Tate--Hochschild cohomology,
Transactions of the American Mathematical Society, 374 (2021), 7, 4537--4577.

\bibitem{p3}
Xu P.,
Gerstenhaber algebras and BV-algebras in Poisson geometry,
Communications in Mathematical Physics, 200 (1999), 3, 545--560.

\end{thebibliography}

\end{document}
